\section{Ordered Hurwitz germs and a family closing in ordinary zeta jets}

Fix $a>0$. All logarithms of $n+a$ are real. Write
\[
 Z_d(s_1,\ldots,s_d;a)
 =\sum_{n_1>\cdots>n_d\ge0}\prod_{j=1}^d(n_j+a)^{-s_j},
 \qquad Z_0=1.
\]
The defining series converges normally whenever
$\Re(s_1+\cdots+s_j)>j$ for $1\le j\le d$.
Put $u_j=s_j-1$ and $P_k=u_1+\cdots+u_k$.
Empty products equal one.

\subsection{An elementary holomorphic remainder construction}

Define the tail defect
\begin{equation}\label{ord:taildefect}
 D(s,x)=\zeta(s,x+1)-\frac{x^{1-s}}{s-1},\qquad x>0.
\end{equation}
Its apparent singularity at $s=1$ is removable. On $\Re s>0$ it has
the absolutely convergent representation
\begin{equation}\label{ord:tailmellin}
 D(s,x)=\frac1{\Gamma(s)}\int_0^\infty
 t^{s-1}e^{-xt}\left(\frac1{e^t-1}-\frac1t\right)\,dt.
\end{equation}
Initially this follows by subtracting the usual Mellin integrals in
$\Re s>1$. The bracket is bounded near zero and is $O(1/t)$ at
infinity, proving continuation and normal convergence in $\Re s>0$.
In particular, uniformly for $s$ on a compact subset there,
$D(s,x)=O(x^{-\Re s})$ as $x\to\infty$.

For $d\ge2$, define
\begin{equation}\label{ord:Edef}
 E_d(u_1,\ldots,u_d;a)
 =\sum_{n_2>\cdots>n_d\ge0}
 D(1+u_1,n_2+a)\prod_{j=2}^d(n_j+a)^{-1-u_j}.
\end{equation}
This is holomorphic on the $\ell^1$ ball
\[
 \mathcal U_d=\{\mathbf u\in\mathbb C^d:
                         |u_1|+\cdots+|u_d|<1/2\}.
\]
Indeed, on a compact subset the inner finite sum over $n_3,\ldots,n_d$
is bounded by a constant times
$(1+n_2)^{\sum_{j=3}^d|u_j|}(1+\log(1+n_2))^{d-2}$.
The remaining outer summand is therefore
$O((1+n_2)^{-2+\sum|u_j|}(1+\log(1+n_2))^{d-2})$.
This is summable uniformly on compact subsets. Holomorphy follows by
normal convergence; fixed spectral derivatives can also be taken termwise.

Set $R_0=1$ and
\[
 R_1(u;a)=\zeta(1+u,a)-\frac1u.
\]
Recursively, for $d\ge2$, put
\begin{equation}\label{ord:Rrecursion}
 R_d(u_1,\ldots,u_d;a)=E_d(u_1,\ldots,u_d;a)
 +\frac{R_{d-1}(u_1+u_2,u_3,\ldots,u_d;a)
             -R_{d-1}(u_2,u_3,\ldots,u_d;a)}{u_1}.
\end{equation}
The last quotient is its holomorphic divided difference at $u_1=0$;
equivalently, it equals the integral from zero to one of the first
partial derivative of $R_{d-1}$ at
$(u_2+\tau u_1,u_3,\ldots,u_d)$. Thus every $R_d$ is holomorphic on
$\mathcal U_d$.

\begin{theorem}[Canonical shifted ordered decomposition]
\label{ord:canonical}
As meromorphic germs on $\mathcal U_d$,
\begin{equation}\label{ord:canonicalformula}
 \boxed{\displaystyle
 Z_d(1+u_1,\ldots,1+u_d;a)
 =\sum_{k=0}^{d}
 \frac{R_{d-k}(u_{k+1},\ldots,u_d;a)}{P_1P_2\cdots P_k}.}
\end{equation}
The local polar hyperplanes are precisely $P_k=0$, each generically
simple in the independent prefix coordinates.
\end{theorem}

\begin{proof}
Summing first over $n_1>n_2$ in the absolute-convergence domain gives
\[
 Z_d(1+\mathbf u;a)
 =\frac1{u_1}Z_{d-1}(1+u_1+u_2,1+u_3,\ldots,1+u_d;a)
       +E_d(\mathbf u;a).
\]
Apply the induction hypothesis to the depth $d-1$ function. Every term
except its holomorphic remainder already has the required denominator
$P_1\cdots P_k$. Splitting the remaining numerator into
$R_{d-1}(u_2,\ldots,u_d)$ plus its difference gives exactly
\eqref{ord:Rrecursion}. Normal convergence and the identity theorem then
establish the meromorphic identity throughout $\mathcal U_d$.
In the coordinates $P_1,\ldots,P_d$, every displayed denominator is
square-free. The coefficient of $P_1^{-1}\cdots P_d^{-1}$ is one,
because only the $k=d$ term contributes. Consequently no one of these
polar divisors is removable as a germ.
\end{proof}

For $a=1$, the shape of \eqref{ord:canonicalformula} is the ordered
Laurent decomposition of Saha\cite[Theorem 2]{Saha}. The construction above
provides a short proof for arbitrary common shift $a>0$ and supplies
absolutely convergent representatives for every regular coefficient.
The canonical framework is credited to that work; the explicit shifted
construction and the evaluations below are derived here.

\subsection{Explicit coefficients in harmonic--Stieltjes coordinates}

Let
\[
 Q_m(x)=\gamma_m(x+1)+\frac{\log^{m+1}x}{m+1}.
\]
In particular $Q_0(x)=\log x-\psi(x+1)$, and
\[
 D(1+u,x)=\sum_{m\ge0}\frac{(-1)^m}{m!}Q_m(x)u^m.
\]
Consequently, for a multi-index $\mathbf m=(m_1,\ldots,m_d)$,
\begin{equation}\label{ord:Ecoeff}
 [\mathbf u^{\mathbf m}]E_d(\mathbf u;a)
 =\frac{(-1)^{|\mathbf m|}}{m_1!\cdots m_d!}
 \sum_{n_2>\cdots>n_d\ge0}
 \frac{Q_{m_1}(n_2+a)\log^{m_2}(n_2+a)}{n_2+a}
 \prod_{j=3}^d\frac{\log^{m_j}(n_j+a)}{n_j+a}.
\end{equation}
These series are absolutely convergent. For a fixed outer index they
are finite nested generalized harmonic sums with logarithmic insertions.
If $r^{(d)}_{\mathbf m}=[\mathbf u^{\mathbf m}]R_d$, then
\begin{equation}\label{ord:Rcoeffrec}
 r^{(d)}_{m_1,\ldots,m_d}
 =[\mathbf u^{\mathbf m}]E_d
 +\binom{m_1+m_2+1}{m_1+1}
 r^{(d-1)}_{m_1+m_2+1,m_3,\ldots,m_d}.
\end{equation}
This is a finite-depth, exact coefficient algorithm, with base
$r^{(1)}_m=(-1)^m\gamma_m(a)/m!$.

\subsection{The diagonal regular germ and its Gamma normalization}

Write $r_d(t;a)=R_d(t,\ldots,t;a)$, with $r_0=1$.
Define the generating function, formally in $y$ and analytically for
sufficiently small $y,t$,
\begin{equation}\label{ord:Fdef}
 F_a(y,t)=\exp\left\{
 y\left(\zeta(1+t,a)-\frac1t\right)
 +\sum_{m\ge2}\frac{(-1)^{m-1}}m
                   \zeta(m+mt,a)y^m\right\}.
\end{equation}

\begin{theorem}[Diagonal regular coefficients]
\label{ord:diag}
For every $d\ge0$,
\begin{equation}\label{ord:diagformula}
 r_d(t;a)=[y^d]F_a(y,t),\qquad
 \sum_{d\ge0}r_d(0;a)y^d=\frac{\Gamma(a)}{\Gamma(a+y)}.
\end{equation}
\end{theorem}

\begin{proof}
For $\Re t>0$, the product over the common lattice gives the standard
elementary-symmetric generating series
\[
 \sum_{d\ge0}Z_d(1+t,\ldots,1+t;a)y^d
 =\prod_{n\ge0}\left(1+\frac{y}{(n+a)^{1+t}}\right)
 =\exp\sum_{m\ge1}\frac{(-1)^{m-1}}m\zeta(m+mt,a)y^m.
\]
On the diagonal, \eqref{ord:canonicalformula} becomes
$Z_d=\sum_{k=0}^dr_{d-k}(t)/(k!t^k)$.
Thus multiplication by $\exp(-y/t)$ gives \eqref{ord:Fdef},
coefficient by coefficient. Holomorphic continuation reaches $t=0$.
Finally, the Taylor series of $\log\Gamma(a+y)$ is
\[
 \log\Gamma(a)+\psi(a)y
       +\sum_{m\ge2}\frac{(-1)^m}{m}\zeta(m,a)y^m.
\]
Since $\gamma_0(a)=-\psi(a)$, exponentiation proves the last assertion.
\end{proof}

Let $B_d(a)=[y^d]\Gamma(a)/\Gamma(a+y)$.
These are explicit Bell polynomials in digamma and polygamma values.
For example, abbreviating $g=\gamma_0(a)$ and $z_m=\zeta(m,a)$,
\begin{align*}
 B_1&=g,& B_2&=\frac{g^2-z_2}{2},\\
 B_3&=\frac{g^3-3gz_2+2z_3}{6},\\
 B_4&=\frac{g^4-6g^2z_2+3z_2^2+8gz_3-6z_4}{24},\\
 B_5&=\frac{g^5-10g^3z_2+15gz_2^2+20g^2z_3
                   -20z_2z_3-30gz_4+24z_5}{120}.
\end{align*}

\subsection{All depths with one distinguished slope}

The finite part below means the coefficient of $\varepsilon^0$ after
the stated one-variable substitution. It is not a value assigned at the
intersection of polar divisors.

\begin{theorem}[A continuum of exact ordered finite parts at every depth]
\label{ord:oneslope}
Let $d\ge1$, and let $p,q\in\mathbb C$ satisfy
$p+jq\ne0$ for $0\le j\le d-1$. Put
$D_k(p,q)=\prod_{j=0}^{k-1}(p+jq)$.
Then
\begin{align}\label{ord:oneslopeformula}
 &\operatorname{FP}_{\varepsilon=0}
 Z_d(1+p\varepsilon,1+q\varepsilon,\ldots,1+q\varepsilon;a)
 \notag\\
 &\qquad=\boxed{\displaystyle B_d(a)+
 \sum_{k=1}^{d-1}\frac{q^k}{D_k(p,q)}
                      [y^{d-k}t^k]F_a(y,t).}
\end{align}
Every expression on the right is a finite polynomial in ordinary
Hurwitz zeta derivatives and generalized Stieltjes constants, with
coefficients rational in $p,q$. It uses only $\gamma_j(a)$ with
$0\le j\le d-1$ and $\zeta^{(j)}(m,a)$ with
$m\ge2$, $j\ge0$, and $m+j\le d$.
The case $q=0$ is included and gives simply $B_d(a)$.
\end{theorem}

\begin{proof}
For $k\ge1$, the tail in \eqref{ord:canonicalformula} is diagonal,
and its denominator becomes $D_k(p,q)\varepsilon^k$.
Its constant coefficient is therefore
$q^k[t^k]r_{d-k}(t;a)/D_k(p,q)$. The $k=0$ term contributes
$R_d(0)=B_d$, and the $k=d$ term is a pure pole and contributes zero.
The coefficient formula \eqref{ord:Fdef} proves the closure and the
stated finite range of orders.
\end{proof}

The same argument evaluates the complete polar part. More precisely,
for $-d\le\ell\le0$,
\begin{equation}\label{ord:allnonpositive}
 [\varepsilon^\ell]Z_d(1+p\varepsilon,1+q\varepsilon,\ldots;a)
 =\sum_{k=\max(0,-\ell)}^d
 \frac{q^{k+\ell}}{D_k(p,q)}
 [y^{d-k}t^{k+\ell}]F_a(y,t),
\end{equation}
where $D_0=1$ and $0^0=1$. For $k=0$ this formula is used only
when $\ell=0$, where its contribution is $B_d$.
Denote the finite part in \eqref{ord:oneslopeformula} by
$\mathcal F_d(p,q;a)$. For $q\ne0$ it depends only on $\lambda=p/q$:
\begin{equation}\label{ord:factorialdirection}
 \mathcal F_d(p,q;a)=B_d(a)+
 \sum_{k=1}^{d-1}\frac{[y^{d-k}t^k]F_a(y,t)}{(\lambda)_k}.
\end{equation}
Thus its dependence on the distinguished slope ratio is an explicit
finite inverse-factorial expansion. Formula \eqref{ord:factorialdirection}
can have a meromorphic continuation in $\lambda$ beyond the transverse
directions; such a continuation of a coefficient must not be confused
with restriction of the original multiple-zeta germ to a polar divisor.

This is an exact evaluation theorem, not a claim of algebraic
independence or minimality of the resulting constants. The general
ordered problem is more complicated: independent trailing slopes can
retain nonsymmetric multiple Stieltjes coefficients.

\subsection{Explicit depths two through five}

Set $g_j=\gamma_j(a)$, $g=g_0$, and
$z_m^{[j]}=\partial_s^j\zeta(s,a)|_{s=m}$; thus
$z_m=z_m^{[0]}$. Then
\begin{align}\label{ord:F2}
 \mathcal F_2&=B_2-\frac qp g_1,\\
 \mathcal F_3&=B_3-\frac qp(gg_1+z_2^{[1]})
                  +\frac{q^2g_2}{2p(p+q)}.\label{ord:F3}
\end{align}
At depth four one obtains
\begin{align}\label{ord:F4}
 \mathcal F_4={}&B_4+
 \frac qp\left(-\frac{g^2g_1}{2}+\frac{g_1z_2}{2}
                        -gz_2^{[1]}+z_3^{[1]}\right)\notag\\
 &+\frac{q^2}{p(p+q)}
       \left(\frac{g_1^2+gg_2}{2}-z_2^{[2]}\right)
 -\frac{q^3g_3}{6p(p+q)(p+2q)}.
\end{align}
At depth five,
\begin{align}\label{ord:F5}
 \mathcal F_5={}&B_5+\frac qp A_{4,1}
 +\frac{q^2}{p(p+q)}A_{3,2}
 +\frac{q^3}{p(p+q)(p+2q)}A_{2,3}
 +\frac{q^4g_4}{24p(p+q)(p+2q)(p+3q)},
\end{align}
where
\begin{align*}
 A_{4,1}={}&-\frac{g^3g_1}{6}+\frac{gg_1z_2}{2}
 -\frac{g^2z_2^{[1]}}2+\frac{z_2z_2^{[1]}}2
 -\frac{g_1z_3}{3}+gz_3^{[1]}-z_4^{[1]},\\
 A_{3,2}={}&\frac{gg_1^2}{2}+\frac{g^2g_2}{4}
 -\frac{g_2z_2}{4}+g_1z_2^{[1]}-gz_2^{[2]}
 +\frac32z_3^{[2]},\\
 A_{2,3}={}&-\frac{gg_3}{6}-\frac{g_1g_2}{2}
 -\frac23z_2^{[3]}.
\end{align*}
These formulas can be checked exactly using Newton's identities, with
independent formal symbols for every Hurwitz jet.

\subsection{Arbitrary depth-three directions and the one extra constant}

Define
\begin{equation}\label{ord:etadef}
 \eta(a)=\frac{\gamma_2(a)}2-
 \sum_{n\ge0}\frac{\gamma_1(n+a+1)+\tfrac12\log^2(n+a)}{n+a}.
\end{equation}
The summand is $O((1+\log n)/n^2)$.
The identities above imply
\[
 \eta(a)=[u]R_2(u,0;a),\qquad
 [v]R_2(0,v;a)=-\gamma_0(a)\gamma_1(a)-\zeta'(2,a)-\eta(a).
\]
For the second identity use the depth-two stuffle law, which gives
\[
 R_2(u,v;a)+R_2(v,u;a)
 =R_1(u;a)R_1(v;a)-\zeta(2+u+v,a).
\]

\begin{theorem}[Every transverse straight direction at depth three]
\label{ord:threedirections}
Let $c_1,c_2,c_3\in\mathbb C$ with
$c_1(c_1+c_2)(c_1+c_2+c_3)\ne0$. Then
\begin{align}\label{ord:threeFP}
 &\operatorname{FP}_{\varepsilon=0}
 Z_3(1+c_1\varepsilon,1+c_2\varepsilon,1+c_3\varepsilon;a)
 \notag\\
 &\quad=B_3(a)-\frac{c_3}{c_1}
       \bigl(\gamma_0(a)\gamma_1(a)+\zeta'(2,a)\bigr)
 +\frac{c_3^2\gamma_2(a)}{2c_1(c_1+c_2)}
 +\frac{c_2-c_3}{c_1}\eta(a).
\end{align}
The complete polar part is
\begin{align}\label{ord:threepoles}
 &\frac{1}{c_1(c_1+c_2)(c_1+c_2+c_3)\varepsilon^3}
 +\frac{\gamma_0(a)}{c_1(c_1+c_2)\varepsilon^2}\notag\\
 &\quad+\frac1\varepsilon
 \left\{\frac{B_2(a)}{c_1}
       -\frac{c_3\gamma_1(a)}{c_1(c_1+c_2)}\right\}.
\end{align}
\end{theorem}

\begin{proof}
At depth three the canonical identity reads
\[
 Z_3=\frac1{u_1(u_1+u_2)(u_1+u_2+u_3)}
 +\frac{R_1(u_3)}{u_1(u_1+u_2)}
 +\frac{R_2(u_2,u_3)}{u_1}+R_3(u_1,u_2,u_3).
\]
Expand $R_1$ through quadratic order and $R_2$ through linear order;
use $R_2(0,0)=B_2$, $R_3(0,0,0)=B_3$, and the two displayed
first coefficients of $R_2$. This gives both assertions.
\end{proof}

The coefficient of $\eta$ vanishes exactly on the family $c_2=c_3$
(without any assertion that $\eta$ is independent of the other constants,
or that it is nonzero at every shift).
The full labelled symmetrization also cancels it, because the two
orders of the last two slopes have opposite coefficients.
Thus generic depth-three directional finite parts need only one
nonsymmetric depth-two coefficient beyond ordinary Hurwitz data;
no new depth-three constant is required for the constant Laurent term.

